\chapter{Solid State Physics: Midterm Exam, April 30, 2015}
\label{q:physics-2015-04-midterm}
\begingroup
\setlength{\parskip}{0.35em}
\setlist[enumerate]{itemsep=0.25em,topsep=0.3em}
\textbf{Printed total: 100\%.}
\begin{sourceNote}
Reconstructed from two pages of the source exam, retaining question numbers, marks, and English question content. Handwritten annotations are excluded from the question text.
\end{sourceNote}

\section*{1. True or false (15\%)}
True or false. Explain your answer.
Crystalline BaTiO$_3$ is underlain by a simple cubic lattice with a basis consisting of five atoms: a Ba atom, a Ti, and three O atoms.
\begin{enumerate}[label=(\alph*)]
\item (5\%) A vector connecting two Ti atoms must be a translation vector of the lattice.
\item (5\%) A vector connecting a Ba atom with a Ti atom must not be a translation vector of the lattice.
\item (5\%) A vector connecting two O atoms must be a translation vector of the lattice.
\end{enumerate}

\section*{2. Direct and reciprocal lattices (30\%)}
Consider a Bravais lattice spanned by the primitive vectors
\[
\vct t_1=\frac a2(0,1,1),\qquad
\vct t_2=\frac a2(1,0,1),\qquad
\vct t_3=\frac a2(1,1,0).
\]
\begin{enumerate}[label=(\alph*)]
\item (10\%) Find the volumes of the Wigner--Seitz cell and the first Brillouin zone.
\item (5\%) Find the three primitive vectors $\vct g_1,\vct g_2,\vct g_3$ of the reciprocal lattice, satisfying
$\vct t_i\cdot\vct g_j=2\pi\delta_{ij}$, $i,j=1,2,3$.
\item (5\%) Consider a family of direct lattice planes that is normal to $\vct t_1$. Find the distance between two successive lattice planes.
\item (5\%) The lattice points on each lattice plane normal to $\vct t_1$ form a two-dimensional lattice. Find the density of lattice points (i.e., the number of lattice points per unit area) on the two-dimensional lattice plane.
\item (5\%) Can the three vectors
\[
\vct G_1=\frac\pi a(0,1,1),\qquad
\vct G_2=\frac\pi a(1,0,1),\qquad
\vct G_3=\frac\pi a(1,1,0)
\]
span the reciprocal lattice? Explain your answer.
\end{enumerate}

\section*{3. Diamond structure (5\%)}
A crystal in diamond structure has lattice constant $a$. Find the density of atoms (i.e., the number of atoms per unit volume).

\section*{4. Definitions (20\%)}
Describe the following terms:
\begin{enumerate}[label=(\alph*)]
\item (5\%) three-dimensional hexagonal lattice;
\item (5\%) Bloch theorem;
\item (5\%) critical point;
\item (5\%) Fermi energy.
\end{enumerate}

\section*{5. Critical point and density of states (10\%)}
Consider a three-dimensional crystal with a critical point $\vct k_1$ of type $M_1$.
\begin{enumerate}[label=(\alph*)]
\item (5\%) Express the energy dispersion for $\vct k$ in a small neighborhood of the critical point.
\item (5\%) Find the density of states (as a function of energy) for the states in the neighborhood.
\end{enumerate}

\section*{6. Density parameter of a free-electron gas (10\%)}
Consider a three-dimensional free-electron gas with a dimensionless parameter $r_s$ for its electron density. Express in terms of $r_s$ the electron density $n$, the Fermi wavenumber $k_F$, the Fermi energy $\varepsilon_F$, the Fermi temperature $T_F$, and the average electron energy for the gas at the ground state.

\section*{7. Electronic heat capacity (5\%)}
Consider a three-dimensional gas of $N$ free electrons. Explain that the gas has a heat capacity
\[
C_V\approx\frac32\kb N\frac{T}{T_F}
\]
rather than the classical $C_V=\frac32\kb N$.
\begin{sourceNote}
The estimate is retained as printed. The solutions explain why its numerical coefficient differs from the rigorous low-temperature expansion.
\end{sourceNote}

\section*{8. Fermi--Dirac tail (5\%)}
Explain why the Fermi--Dirac distribution function in a simple metal can be approximated well as an exponential function for electrons with energy higher than the vacuum level.
\endgroup
